\section{Анализ инструментов профилирования Go и подходов к интеграции профилирования в CI/CD}
\label{sec:razdel1}

\subsection{Обзор инструментов сбора и анализа данных о производительности в экосистеме Go}
\label{sec:1.1}

Язык Go~\cite{go-spec} изначально проектировался с расчётом на встроенную
поддержку измерения производительности: инструментарий профилирования и
бенчмаркинга поставляется вместе со стандартной библиотекой и
инструментами \texttt{go}, а не выносится в сторонние пакеты, как это
исторически происходило, например, в экосистеме Java (JMH подключается
отдельно) или Node.js (clinic.js, 0x)~\cite{donovan-kernighan}.

\textbf{Бенчмарки (\texttt{testing.B}).} Пакет \texttt{testing} предоставляет тип
\texttt{testing.B}, позволяющий описывать микробенчмарки функцией вида
\texttt{func BenchmarkXxx(b *testing.B)}, где тело цикла выполняется \texttt{b.N}
раз — количество итераций подбирается средой выполнения автоматически, пока не
будет достигнута статистически достаточная длительность измерения. Команда
\texttt{go test -bench=. -benchmem} выполняет такие функции и печатает три базовые
метрики на итерацию: время (\texttt{ns/\allowbreak op}), объём выделенной памяти
(\texttt{B/\allowbreak op}) и число аллокаций (\texttt{allocs/\allowbreak op})~\cite{go-testing-doc}.
Многократный запуск с флагом \texttt{-count=N} даёт выборку из $N$ независимых
измерений на каждый бенчмарк, что необходимо для последующего статистического
сравнения, поскольку однократное измерение производительности крайне
чувствительно к шуму окружения (частота планировщика ОС, состояние кэшей
процессора, фоновая нагрузка CI-раннера)~\cite{georges2007,mytkowicz2009}.

\textbf{\texttt{benchstat}.} Утилита \texttt{golang.org/\allowbreak x/\allowbreak perf/\allowbreak cmd/\allowbreak benchstat}
принимает один или два текстовых отчёта \texttt{go test -bench} и вычисляет
для каждого бенчмарка медиану, доверительный интервал и, при сравнении двух
наборов, относительное изменение метрики с $p$-значением критерия
Манна --- Уитни (по умолчанию) или критерия Уэлча, если данные могут считаться
нормально распределёнными~\cite{benchstat-doc,welch1947}. Именно
\texttt{benchstat} лежит в основе сравнения "текущий прогон vs. эталон" в
данной работе (раздел~\ref{sec:razdel4}) — однако сам по себе он не принимает
решения "регрессия / не регрессия": он лишь поставляет числа, интерпретация
которых (выбор порога, уровня значимости) остаётся на усмотрение вызывающей
стороны.

\textbf{\texttt{pprof} (\texttt{runtime/\allowbreak pprof}, \texttt{net/\allowbreak http/\allowbreak pprof}).}
В отличие от агрегированных чисел бенчмарка, пакет \texttt{runtime/\allowbreak pprof}
позволяет получить профиль — набор сэмплов стека вызовов, размеченных по
времени (CPU-профиль, сэмплирование по таймеру) или по объёму выделенной
памяти (heap/alloc-профиль, сэмплирование аллокатора). Профиль сохраняется в
бинарном формате \texttt{pprof} (сжатый protobuf) и анализируется утилитой
\texttt{go tool pprof}, которая умеет строить текстовые топ-листы
(\texttt{-top}), графы вызовов (\texttt{-web}), построчные разборы горячих
функций (\texttt{-list}) и, что ключево для настоящей работы, \emph{разность
двух профилей} через флаг \texttt{-diff\_base}, вычитающую сэмплы одного
профиля из другого функция за функцией~\cite{go-pprof-doc,go-pprof-tool,pprof-github}.
Команда \texttt{go test} умеет снимать такие профили прямо во время прогона
бенчмарков флагами \texttt{-cpuprofile} и \texttt{-memprofile} — это
единственный источник данных, показывающий, \emph{какая именно функция}
стала "тяжелее", тогда как \texttt{benchstat} видит только агрегат по
бенчмарку в целом.

\textbf{\texttt{go tool trace}.} Отдельный инструмент, визуализирующий
исполнение горутин, работу планировщика и сборщика мусора во времени;
полезен для диагностики задержек, вызванных конкуренцией за ресурсы
(блокировки, GC-паузы), но не рассматривается подробно в данной работе,
поскольку тестовое приложение (раздел~\ref{sec:razdel3}) однопоточное по
своей логике и не порождает конкурентного поведения, которое требовало бы
такого анализа.

\textbf{Промышленные системы непрерывного профилирования.} Отдельный класс
инструментов — системы, которые снимают pprof-совместимые (или
eBPF-based) профили \emph{непрерывно в продакшене}, а не разово в CI:
инфраструктура Google-Wide Profiling, впервые описанная
в~\cite{ren2010}, стала прообразом современных open-source решений —
Grafana Pyroscope~\cite{pyroscope-doc} и Parca (Polar Signals)~\cite{parca-doc},
а также коммерческих сервисов (Google Cloud Profiler, Datadog Continuous
Profiler). Эти системы решают принципиально другую задачу: они дают
постоянный поток профилей с боевых инстансов для ретроспективного анализа
("что происходило в проде в момент инцидента"), а не блокируют слияние кода
с обнаруженной регрессией. Практика эксплуатационного профилирования тесно
связана с более широкой дисциплиной сопровождения production-систем
(site reliability engineering), где постоянный сбор телеметрии, включая
профили производительности, используется прежде всего для диагностики уже
проявившихся инцидентов, а не для их предотвращения на этапе разработки
кода~\cite{beyer-sre}. Разграничение этих двух смыслов термина
"непрерывное профилирование" — эксплуатационный (production continuous
profiling) и процессный (профилирование при каждом прогоне CI до слияния) —
формализуется в конце данного раздела и используется как рабочее определение
на протяжении всей работы.

\subsection{Сравнительный анализ подходов к интеграции профилирования и обнаружения регрессий в CI/CD}
\label{sec:1.2}

Обнаружение регрессий производительности до попадания кода в основную ветку
концептуально относится к тому же классу задач, что и статический анализ или
модульное тестирование в CI, но на практике встречается в пайплайнах заметно
реже: по данным обзора существующих open-source решений (таблица~\ref{tab:1.1})
подавляющее большинство CI-конфигураций Go-проектов ограничивается
\texttt{go vet}, \texttt{go test} и линтерами, не включая шаг сравнения
производительности вовсе — go test сам по себе не проваливает сборку при
деградации метрик, а лишь сообщает средние значения.

\begin{table}[htbp]
\caption{Сравнение подходов к обнаружению регрессий производительности в CI/CD}
\label{tab:1.1}
\centering
\begin{tabular}{p{2.6cm}p{3.6cm}p{3.6cm}p{4.2cm}}
\toprule
Подход & Уровень данных & Статистическая строгость & Ограничение \\
\midrule
Ручное профилирование разработчиком & pprof-профиль, разово & Отсутствует (субъективная оценка) & Не встроено в пайплайн, зависит от дисциплины разработчика \\
\texttt{go test -bench} без сравнения & Агрегат бенчмарка, разово & Отсутствует & Числа печатаются, но не с чем сравнивать автоматически \\
\texttt{benchstat} в CI (Bencher.dev, CodSpeed, GitHub Action \texttt{benchmark-action}) & Агрегат бенчмарка, серия & Есть ($p$-значение, доверительный интервал) & Не показывает, какая функция деградировала — только итоговое число \\
JMH + CI-визуализация (экосистема Java) & Агрегат бенчмарка, серия & Есть (аналогично benchstat) & Не Go-специфично, приводится как референс зрелости подхода в другой экосистеме \\
Production continuous profiling (Pyroscope, Parca, Google-Wide Profiling) & pprof/eBPF-профиль, непрерывно & Не применяется (не задача обнаружения регрессий) & Работает постфактум, после выкладки в прод, не блокирует PR \\
\textbf{Прототип данной работы} & Агрегат (benchstat) \emph{и} pprof-профиль, при каждом PR & Есть (порог + $p$-значение) & Синтетическая нагрузка бенчмарка может не отражать прод-трафик \\
\bottomrule
\end{tabular}
\end{table}

Из таблицы~\ref{tab:1.1} видно, что существующие решения занимают две
крайние точки: либо статистически строгое, но "слепое" к причине сравнение
агрегатов (benchstat-класс инструментов), либо детальный, но не
автоматизированный и не привязанный к порогу профиль (pprof-класс
инструментов и системы непрерывного профилирования). Ни один из
рассмотренных инструментов не совмещает оба уровня данных в одном шаге
CI-пайплайна. Это наблюдение определяет техническое решение,
принятое в разделах~\ref{sec:razdel3}--\ref{sec:razdel4}: прототип запускает
на каждый pull request как \texttt{benchstat}-сравнение (решение "регрессия/
не регрессия"), так и параллельный съём pprof-профилей с их
\texttt{diff\_base}-сравнением (объяснение \emph{какая функция} стала
причиной), сохраняя оба артефакта.

Дополнительно был проведён обзор литературы по статистической методологии
сравнения производительности программ. Работы Georges et al.~\cite{georges2007}
и Mytkowicz et al.~\cite{mytkowicz2009} независимо показывают, что
единичные измерения времени выполнения программ систематически вводят в
заблуждение из-за детерминированных, но невидимых наблюдателю эффектов
окружения (выравнивание кода и данных в памяти, порядок переменных
окружения ОС), и что корректное сравнение обязано опираться на выборки
и статистические критерии, а не на сравнение единственных чисел "было/стало".
Это напрямую обосновывает использование в данной работе критерия Уэлча
(либо непараметрического аналога) поверх выборки из \texttt{-count} повторов,
а не сравнение единичных прогонов.

\subsection{Выводы}

\begin{enumerate}
\item Экосистема Go предоставляет развитый инструментарий сбора данных о
производительности (\texttt{testing.B}, \texttt{runtime/\allowbreak pprof},
\texttt{go tool pprof}, \texttt{golang.org/\allowbreak x/\allowbreak perf/\allowbreak cmd/\allowbreak benchstat}), но не
поставляет готового решения, принимающего бинарное решение "регрессия
обнаружена / не обнаружена" и блокирующего слияние кода на его основе.
\item Сравнительный анализ (раздел~\ref{sec:1.2}, таблица~\ref{tab:1.1})
показал, что существующие подходы делятся на статистически строгие, но
агрегированные (benchstat-класс) и детальные, но не автоматизированные
относительно порога (pprof-класс, включая промышленные системы
непрерывного профилирования продакшен-систем — Pyroscope, Parca,
Google-Wide Profiling); совмещения обоих уровней данных в одном шаге CI не
обнаружено.
\item Показано (Georges et al.~\cite{georges2007}, Mytkowicz et
al.~\cite{mytkowicz2009}), что сравнение единичных измерений
производительности статистически некорректно из-за скрытых эффектов
окружения выполнения — критерий обнаружения регрессии обязан опираться на
выборку повторных измерений и статистический критерий значимости, а не на
разность двух отдельных чисел.
\item Определены два уровня метрик, необходимых для полноценного
обнаружения и диагностики регрессии: агрегированные метрики бенчмарка
(\texttt{ns/\allowbreak op}, \texttt{B/\allowbreak op}, \texttt{allocs/\allowbreak op}) и функциональные метрики
профиля (\texttt{flat}/\texttt{cumulative} время CPU-профиля,
\texttt{alloc\_space} heap-профиля).
\item Зафиксировано рабочее разграничение термина "непрерывное
профилирование": в данной работе оно означает профилирование, встроенное в
каждый прогон CI-пайплайна (\emph{процессный} смысл), в отличие от
постоянного профилирования уже выложенного в продакшен приложения
(\emph{эксплуатационный} смысл, реализуемый системами вроде Pyroscope/Parca)
— накладные расходы именно процессного профилирования количественно не
изучены в рассмотренных источниках применительно к Go-бенчмаркам, что
обосновывает необходимость эксперимента раздела~\ref{sec:5.3}.
\end{enumerate}

\subsection{Цели и задачи УИР/НИР}

\textbf{Цель работы} — разработать и экспериментально апробировать прототип
системы, интегрирующей два уровня данных о производительности
Go-приложения — статистическое сравнение бенчмарков (\texttt{benchstat}) и
pprof-диагностику на уровне функций — в CI/CD-пайплайн, чтобы регрессии
производительности обнаруживались автоматически на этапе pull request, до
попадания кода в основную ветку, а не постфактум в эксплуатации.

По результатам анализа, проведённого в разделе~\ref{sec:razdel1}, были
уточнены следующие задачи работы (аналитические задачи, выполнявшиеся на
этапе РСПЗ, здесь заменены указанием, как их результат повлиял на решение
остальных задач):

\begin{enumerate}
\item Результаты анализа (раздел~\ref{sec:1.2}) показали отсутствие
решения, совмещающего агрегированное сравнение и pprof-диагностику, —
это определило требование к архитектуре прототипа (раздел~\ref{sec:razdel3})
включить оба механизма как параллельные, а не взаимоисключающие шаги
пайплайна.
\item Формализовать критерий обнаружения регрессии производительности,
опирающийся на статистическую значимость и практическую величину отклонения
одновременно, а также модель процесса профилирования в CI/CD-пайплайне
(раздел~\ref{sec:razdel2}).
\item Спроектировать архитектуру прототипа: тестовое Go-приложение с
воспроизводимой нагрузкой для бенчмарков, формат хранения эталонных данных
(агрегированного и pprof-профилей), модуль сравнения и CI-пайплайн
(раздел~\ref{sec:razdel3}).
\item Реализовать прототип в виде исполняемых артефактов — тестовое
приложение с бенчмарками, скрипты сравнения с эталоном (\texttt{benchstat})
и pprof-диагностики (\texttt{diff\_base}), CI-пайплайн на GitHub Actions
(раздел~\ref{sec:razdel4}).
\item Провести экспериментальную апробацию на контролируемых сценариях
регрессий: оценить корректность обнаружения (долю ложных срабатываний и
пропусков) и количественно оценить накладные расходы, которые
профилирование добавляет к длительности CI-джобы (раздел~\ref{sec:razdel5}).
\end{enumerate}
